% !TEX root = userguide.tex

\chapter{Factor $\tau$ and $\text{Pe}_h$ in the SUPG method}
\pgst{tfemchA}
\label{chap:tausupg}

\def\chaptertitle{Factor $\tau$ and $\text{Pe}_h$ in the SUPG method}
\def\headdate{11th February 2020}

The stabilization factor $\tau$ in the SUPG
method is given by (see Eq.~\eqref{eq:taufactor}):
\begin{equation}
       \tau = \frac{\beta h}{2 U}
    \label{eq:taufactor_A}
\end{equation}
where $\beta$ is a dimensionless parameter,
$h$ is a typical size of the element in the direction of velocity and $U$ is a
characteristic velocity magnitude. For convection-diffusion problems the parameter $\beta$ is taken to be a function of the mesh P\'{e}clet number $\text{Pe}_h$, given by
\begin{equation}
  \mathrm{Pe}_h = \frac{U h}{2\nu}\label{eq:PehA}
\end{equation}
where $\nu$ is a diffusion coefficient (dimension: [L$^2$/T]).
In this appendix the options in \tfem\ for choosing $h$ and $U$ in the computing of $\tau$ and $\mathrm{Pe}_h$ will be discussed.

\section{The characteristic length $h$}

The characteristic size of an element in the direction of the velocity,
 $h$ needs to be chosen carefully.
If taken too small, wiggly solutions are generated, even in
smooth flows like a flow around a sphere. If taken too large, numerical
instabilities can be generated leading to fast blow up of the solution. Furthermore,
for deformed/stretched elements the optimal value of $h$ depends on the direction
of the local velocity vector.

In \tfem\ the choice for $h$ is determined by the parameters \texttt{htype} and 
\texttt{hlocation}, where the latter only indicates whether $h$ is a constant
in the element (\texttt{hlocation=1}; computed in the center of the element) or
 different for each integration point (\texttt{hlocation=2}).  

\subsection{Projection (\texttt{htype=1})}

A possible definition for $h$ is the projection method as suggested by Brooks
and Hughes \cite{Brooks82}:
\begin{equation}
   h = \frac{|\vek h_1\cdot \vek u| + |\vek h_2\cdot \vek u|}{\|\vek u\|}
\end{equation}
where $\vek u$ is the velocity vector and $\vek h_1$ and $\vek h_2$ are defined
in Figure~\ref{fig:appA} for a quadrilateral element.
\begin{figure}[htbp]
\begin{center}
\includegraphics[scale=0.8]{quad}
\caption{Projection in quad element.}
\label{fig:appA}
\end{center}
\end{figure}

It is possible to define one $h$ value for an element and then the velocity in the center of the element is used for defining $h$ (\texttt{hlocation=1}).
It is however possible to define the value $h$ in each integration point separately using the
local velocity vector (\texttt{hlocation=2}). 

Although the projection method works quite well in 2D for viscoelastic fluid
flow, preliminary calculations in 3D seem to suggest that this method is not
very stable for the natural extension to hexahedral elements. Also the method
is quite ad-hoc and it is not obvious how to extend it to triangular and
tetrahedral shaped element. Therefore we will discuss a different method based
on mapping.

\subsection{Mapping (\texttt{htype=2})}

For deriving an easy but general formula for $h$ that can be used for each
element type we assume the following:
\begin{enumerate}
   \item For elements that have a shape that can be
         obtained from the reference element by 
         isotropic expansion only, it is sufficient to have an element size
         $h$ independent from the direction of the velocity. For this we will
         take
          \[
               h = h_0 \times \texttt{1D expansion factor along a line}
          \]
          where $h_0$ is the typical size of the reference element\footnote{The reference
          element for quadrilaterals and hexahedra is taken to be the standard square reference element
          with edge length of 2. For a triangle or tetrahedron the standard reference element is insufficient, since
          the edges are not of the same length. Therefore, the regular triangle/tetrahedron
          having the length of all edges equal to 1 will be used. This requires an additional mapping matrix $\mat F$
          from the standard to the regular reference element of
          \[
              \mat F = \begin{pmatrix}
                          1 & \frac12 \\[3pt]
                          0 & \frac12\sqrt{3}
                       \end{pmatrix}
         \qquad\text{and}\qquad
              \mat F = \begin{pmatrix}
                          1 & \frac12 & \frac12 \\[3pt]
                          0 & \frac12\sqrt{3} &\frac16\sqrt{3} \\[3pt]
                          0 & 0 & \frac13\sqrt{6}
                       \end{pmatrix}
          \]
          for a triangle and a tetrahedron, respectively. For a prism only the triangular base needs to be transformed and the mapping matrix $\mat F$ from the standard to the regular reference element becomes
          \[
              \mat F = \begin{pmatrix}
                          1 & \frac12 & 0 \\[3pt]
                          0 & \frac12\sqrt{3} & 0\\[3pt]
                          0 & 0 & 1 
                       \end{pmatrix}
          \]
          }, for which
          we take the average length derived from the area or volume. So we
          get: $h_0=2$ for quadrilaterals and hexahedra, $h_0=\sqrt[4]{3}/2$ for
          triangles, $h_0=(\sqrt{2}/12)^{\frac13}$ for tetrahedra, 
          $h_0=(\sqrt{3}/2)^{\frac13}$ for prisms and $h_0=\sqrt[3]{4}$ for pyramids.
   \item For elements of arbitrary shape we take a circle (2D) or sphere (3D)
         in the reference element having a diameter of $h_0$. After mapping 
         the circle/sphere onto the actual element we get an ellipse (2D) or
         ellipsoid (3D). We define
         $h$ as the size of the ellipse/ellipsoid (3D) in the
         direction of the velocity. If we denote the
         mapping tensor of the reference element onto the actual element
         by $\ten F$, we find that
         \[
            \frac{\vek u}{\|\vek u\|} h= \ten F\cdot \vek e_0h_0 
         \]
         where $\vek e_0$ is a vector in the reference element on the unit
         circle/sphere. Taking the inverse and the norm we find that
         \[
             h= \frac{h_0\|\vek u\|}{ \|\ten F^{-1}\cdot\vek u\|}
         \]
\end{enumerate}
The idea of using the mapping tensor for obtaining $h$ has been introduced in
\cite{Hughes86}, apparently using a general value of $h_0=2$.

It is possible to define one $h$ value for an element and then the velocity in the center of the element is used for defining $h$ (\texttt{hlocation=1}).
It is however possible to define the value $h$ in each integration point separately using the
local velocity vector (\texttt{hlocation=2}). 

Early calculations seem to suggest that the mapping
technique has (much) better stability properties for 3D viscoelastic flow than
the projection technique.

\subsection{Constant $h$ defined per element, independent from the velocity vector (\texttt{htype=3})}

The most simple method is taking a constant value $h$ for an element only depending on the size of the element and independent of the velocity. For that, we take  $h=\sqrt{A_\text{e}}$ (2D) and $h=\sqrt[3]{V_\text{e}}$ (3D) with
$A_\text{e}$ and $V_\text{e}$ the area and volume of the element, respectively.

NOTE: this choice of $h$ is expected to work best for non-stretched meshes, such as using triangles and tetrahedrons that have close to optimal shape (all angles equal).

\section{The characteristic velocity $U$}

Various possibilities exist for computing a characteristic velocity $U$
in an element:
\begin{enumerate}
   \item The magnitude of the velocity in the center of the element:
\[
         U = \|\vek u_0\|
\]
where $\vek u_0$ is the velocity in the center (of gravity) of the element.
This approach, already suggested in \cite{Brooks82}, works quite well in
regions where the velocity is nearly constant.
However, near the center of a vortex, the value of $U$ can be much smaller than
the velocity magnitude in the integration points leading to a value for $\tau$
that is much too large. This can lead to numerical instabilities.
   \item The average of the magnitude of the velocities in all integration
points:
\[
         U = \frac1N\sum_{i=1}^N\|\vek u_i\|
\]
where $N$ is the number of integration points in the element.
This procedure
works quite well. It resembles the procedure in \cite{Fan97}, where the
average of the values in the vertices is being used. It is claimed in
\cite{Fan97} that this procedure has accuracy advantages in a spectral element
context. It is unclear whether this is true in a standard FEM approach.
   \item The magnitude of the velocity in the integration points:
\[
         U_i = \|\vek u_i\|, \quad i=1,\dots,N
\]
The characteristic velocity is now different for each integration point. 
This procedure also works quite well.
   \item A global velocity scale $U=U_\text{glob}$, the same for all elements.
This procedure has been used in \cite{Lunsmann93}.
\end{enumerate}
